\section{Run the kernel before connecting the file codec}
Use \texttt{packages/034-lossless-video/construction/video-construction/kernel.asm} as the unfinished source. Its small
caller passes five native words to the selected entry and retains the result at
\texttt{kernel\_return}. Edit your own copy cumulatively. The preparation helper
selects data and controls while preserving the entire instruction section.
The ordinary call uses separate output storage; additional observations select exact
left or right output aliases.

For the worked words, the expected XOR values are hexadecimal zero, \texttt{88},
zero, \texttt{10000} and \texttt{FF0000}. The changed-word count is three. The complete
caller prints that count. An unfinished region returns development status four and
prints an under-construction message; invalid argument admission is a different result.
In Debug, stop before the named observation label, inspect its state, then continue.

The prepared input arenas include guard words and use native \texttt{.word} data.
Their serialized memory bytes therefore follow the selected target byte order. Keep
that distinction when comparing them with the big-endian \texttt{MLV1} command counts
and byte-oriented RGB triples. The file codec will expand or pack triples at its own
representation boundaries; the XOR kernel sees native words only.
